\chapter{Results}

This chapter reports the measured output of the local harness. The tables are generated from \texttt{aggregate-result.json}; they must be regenerated after each corpus run.

The latest executed run used run identifier \texttt{20260519-173709-local}. It processed six targets. Four targets built successfully, two targets completed tests successfully, no coverage file was produced, and Docker-backed telemetry execution was unavailable because the qyl compose configuration required a missing \texttt{OPENAI\_API\_KEY} value. These results improve the thesis from a framework-only description to an executed empirical baseline, but they do not yet support a 90+ quality claim.

\section{Corpus}
\begin{tabular}{lll}
\hline
Target & Commit & Dirty status \\
\hline
arqio & 6e831bfad9 & clean \\
otel-conventions-api & 4a67084728 & clean \\
qyl & c2ab2116dd & dirty \\
semconv-testbed & 1a3094e7c2 & clean \\
tourplanner-angular & 8c4941b081 & clean \\
tourplanner-dotnet & 7ae4c644ee & clean \\
\hline
\end{tabular}


\section{Build, Test, and Coverage}
\begin{tabular}{llll}
\hline
Target & Build & Test & Coverage \\
\hline
arqio & failed & skipped & unavailable \\
otel-conventions-api & skipped & skipped & not\_required \\
qyl & passed & failed & unavailable \\
semconv-testbed & passed & passed & not\_required \\
tourplanner-angular & passed & passed & not\_required \\
tourplanner-dotnet & passed & failed & unavailable \\
\hline
\end{tabular}


\begin{tabular}{lll}
\hline
Target & Coverage status & Line coverage \\
\hline
arqio & unavailable & n/a \\
otel-conventions-api & not\_required & n/a \\
qyl & unavailable & n/a \\
semconv-testbed & not\_required & n/a \\
tourplanner-angular & not\_required & n/a \\
tourplanner-dotnet & unavailable & n/a \\
\hline
\end{tabular}


\section{Reproducibility and Classification Agreement}
\begin{tabular}{lll}
\hline
Target & Label & Claim \\
\hline
arqio & service-oriented-or-distributed & agreement only \\
otel-conventions-api & modular-monolith-or-multi-project & agreement only \\
qyl & service-oriented-or-distributed & agreement only \\
semconv-testbed & modular-monolith-or-multi-project & agreement only \\
tourplanner-angular & modular-monolith-or-multi-project & agreement only \\
tourplanner-dotnet & service-oriented-or-distributed & agreement only \\
\hline
\end{tabular}


Ground-truth labels are not configured in the initial corpus. The classification column therefore reports deterministic label agreement and confidence, not correctness. This is a material limitation, but it is still empirical evidence: it shows whether the same harness can process real repositories and produce repeatable structured outputs.

\section{Semantic Convention and Telemetry Evidence}
\begin{tabular}{lll}
\hline
Check & Status & Evidence \\
\hline
qyl stable semantic conventions & passed & build log \\
qyl incubating semantic conventions & passed & build log \\
qyl semantic-conventions package & passed & build log \\
semconv-testbed verification & failed & Docker daemon unavailable \\
runtime telemetry & unavailable & missing compose secret / no spans counted \\
\hline
\end{tabular}


The qyl semantic-convention packages built successfully. The broader semconv-testbed verification did not pass because Docker Desktop was not reachable for its pre-flight check. Runtime telemetry is therefore reported as unavailable for this run rather than inferred from static build success.

\section{Raw Evidence}

The raw evidence is indexed in the appendix. It includes command logs, exit codes, coverage summaries, classification output, and aggregate JSON. Failed or unavailable steps are retained in the result set so that the thesis does not silently exclude inconvenient data.
